\chapter{Von der Kette zur E\texorpdfstring{$_8$}{8}-Naht}
\label{ch:seam}
\section{Ein wichtiger Modellwechsel}
Das Zweizellenmodell benutzt zwei vollständige Vierergraphen mit einer Brücke. Die Vielzellenrechnung in S08 benutzt dagegen Vierersegmente mit nächster Nachbarschaft. Bei $\lambda=J$ wird erst dieses zweite Modell eine uniforme Kette. Derselbe isolierte Singulettzustand garantiert keine identischen Anregungen und keinen identischen Grenzwert.

\begin{figure}[H]\centering
\begin{tikzpicture}[x=1cm,y=1cm]
\foreach \i in {0,...,7}{\fill[navy] (1.4*\i,0) circle (.10);}
\foreach \i in {0,1,2,4,5,6}{\draw[teal,very thick] (1.4*\i,0)--(1.4*\i+1.4,0);}
\draw[amber,very thick] (4.2,0)--node[above,font=\small]{$\lambda$}(5.6,0);
\node[font=\sffamily\small,text=slate] at (2.1,-.65){Vierersegment mit drei Kanten};
\node[font=\sffamily\small,text=slate] at (7.7,-.65){Vierersegment mit drei Kanten};
\node[font=\sffamily\small,text=navy] at (4.9,1){Uniforme Kette erst bei $\lambda=J$};
\end{tikzpicture}
\caption{Die Kettengeometrie der WZW-Vergleichsrechnung. Sie ist weder der vollständige Tetramergraph noch der Clebsch-Graph.}
\end{figure}

S08 berichtet für offene Ketten die folgenden Lücken in Einheiten $J$:
\begin{center}
\begin{tabular}{rrrr}
\toprule $\lambda/J$&8 Träger&12 Träger&16 Träger\\\midrule
$1/4$&0,2664&0,2552&0,2497\\
$1/2$&0,2299&0,2029&0,1890\\
1&0,1759&0,1273&0,1002\\\bottomrule
\end{tabular}
\end{center}
Bei schwacher Brücke ändern sich die endlichen Werte wenig, bei uniformer Kopplung fallen sie. Das stützt einen Phasenunterschied. Es beweist aus diesen wenigen Größen allein keinen strikt positiven thermodynamischen Gap für jedes $\lambda<J$. S06 weist außerdem darauf hin, dass der ursprüngliche 16-Träger-Prüfer eine ausgewählte Darstellungsfamilie untersucht. Eine daraus erhaltene Anregung ist ohne Ausschluss der übrigen Sektoren kein global zertifizierter erster Gap.

\section{Warum die SU(4)-WZW-Theorie hier plausibel ist}
Für den uniformen periodischen Austauschring liefern die Quellen $\Delta n\approx3{,}700$ bei acht und $3{,}663$ bei zwölf Trägern. Der bekannte Vergleich lautet
\begin{equation}
v=\pi/4,\qquad h(4)=3/8,\qquad c=3,
\qquad 2\pi v(h+\bar h)=\frac{3\pi^2}{8}\approx3{,}70110.
\end{equation}
Die Sutherland-Grundenergie pro Träger ist
\begin{equation}
e_\infty=1-\frac{\pi}{8}-\frac34\log2\approx0{,}087439,
\end{equation}
für $H=\sum(1+P)/2$. S08 nennt die Extrapolation $0{,}08773$ und $c_{\mathrm{fit}}=3{,}16$. Diese endlichen Zahlen passen zur $\mathrm{SU}(4)_1$-WZW-Beschreibung des Kettenmodells. Die allgemeine Verbindung von SU(N)-Spinketten und WZW-Niedrigenergiesektoren wird unabhängig in der Literatur untersucht. \href{https://arxiv.org/abs/0806.2563}{Führinger et al. [E5]}

\begin{bild}{Aus vielen Saitenbewegungen entsteht ein glatter Ton}
Ein Gitter besteht aus einzelnen Stellen. Bei großen Wellenlängen kann seine Bewegung durch ein kontinuierliches Feld beschrieben werden, ähnlich wie eine Saite im Alltag glatt wirkt, obwohl Materie nicht kontinuierlich aufgebaut ist. Das bedeutet nicht, dass jeder andere Graph denselben Ton oder dieselbe Feldtheorie besitzt.
\end{bild}

\section{Die Naht braucht mehr als die Zahl acht}
Die zentrale Ladung ist eine Kennzahl einer konformen Feldtheorie. Für die chiralen Niveau-1-Faktoren gilt
\begin{equation}
c(D_5)=\frac{45}{1+8}=5,\quad
c(A_3)=\frac{15}{1+4}=3,\quad
c(E_8)=\frac{248}{1+30}=8.
\end{equation}
Das Produkt hat also schon vor einer Kopplung $c=5+3=8$. Ein später gefundener Wert acht allein ist kein Beweis einer $E_8$-Erweiterung und kein Beweis eines Kondo-Flusses.

Die konkretere Struktur ist die diagonale $\Z_4$-Verklebung:
\begin{center}
\begin{tabular}{ccccc}
\toprule Klasse&$D_5$-Sektor&$A_3$-Sektor&Gewichte&Neue Ströme\\\midrule
0&Vakuum&Vakuum&$0+0$&$45+15$ bei Stufe 1\\
1&16&4&$5/8+3/8=1$&64\\
2&10&6&$1/2+1/2=1$&60\\
3&$\overline{16}$&$\overline4$&$5/8+3/8=1$&64\\\bottomrule
\end{tabular}
\end{center}
Die Erweiterung ergänzt die 60 ursprünglichen Gewicht-1-Ströme um $64+60+64$, insgesamt 248. Auf der Gitterseite entspricht sie den vier Glue-Klassen. S11 nennt
\begin{equation}
\Theta_{E_8}=1+240q+2160q^2+6720q^3+\cdots,
\quad\frac{\Theta_{E_8}}{\eta^8}=q^{-1/3}(1+248q+4124q^2+\cdots).
\end{equation}
Die abstrakte Gittererweiterung ist konkret. Ihr nativer physischer Aufbau muss zusätzlich Operatorprodukte, Cocycle-Phasen, Adjunktion, Domänen und Energieabschätzungen kontrollieren. Der allgemeine Rahmen konformer Einbettungen und Simple-Current-Erweiterungen ist Literaturhintergrund, keine unabhängige TFPT-Bestätigung. \href{https://arxiv.org/abs/1210.6602}{Kac et al. [E6]}

\section{Warum „Z\texorpdfstring{$_4$}{4}-Eichung“ präzisiert werden muss}
S11 ersetzt die frühere Kondo-Zielbeschreibung sinnvoll durch eine konkrete Erweiterungsaufgabe. Die Kurzschreibweise eines Quotienten durch $\Z_4$ kann jedoch verschiedene Operationen meinen. Nur auf invariante Zustände des ursprünglichen Produkts zu projizieren ist nicht dasselbe wie eine Simple-Current-Erweiterung mit den nötigen zusätzlichen beziehungsweise entsprechend gekoppelten Sektoren.

Ein Gittermodell muss angeben, welche Randbedingungen, Ladungssektoren und lokalen Operatoren enthalten sind. Die 248 Zustände auf Stufe eins sind dann ein notwendiger früher Test. Sie allein beweisen noch nicht die vollständigen Operatorprodukte oder den analytischen Skalierungslimes.

S11 schlägt dafür einen SU(4)-Ring zusammen mit zehn Majorana-Moden vor. Deren $D_5$-Sektoren besitzen die Gewichte $0,1/2,5/8,5/8$. Auf der Ringseite sollen die vier Restklassen der Trägerzahl die $A_3$-Sektoren organisieren. Die berichteten endlichen Werte lauten:
\begin{center}
\begin{tabular}{rrrrr}
\toprule $n$&$n\bmod4$&Sektordimension&$E_0/J$&Extrahiertes $x$\\\midrule
5&1&60&0,690983&0,507\\
6&2&180&0,697224&0,460\\
7&3&630&0,634808&0,282\\
8&0&2520&0,539495&$-0,009$\\
9&1&7560&0,898579&0,454\\
10&2&25200&0,998465&0,501\\
11&3&92400&0,993060&0,320\\
12&0&369600&0,944508&$-0,005$\\
13&1&1201200&1,205482&0,431\\\bottomrule
\end{tabular}
\end{center}
Die Werte wurden aus $E_0=ne_\infty-\pi v(c-12x)/(6n)$ gewonnen. In S11 werden die ungeraden Paare gemittelt und mit $3/8$ verglichen, der Zweiersektor mit $1/2$. Diese Sektorzuordnung hängt an den Randbedingungen und endlichen Korrekturen; eine chirale Gewichtstabelle darf nicht unbesehen mit jeder nichtchiralen Ringskalierungsdimension gleichgesetzt werden. Die Daten sind hier als berichtete Hinweise erhalten, nicht neu reproduziert.

\section{Die Chiralitätslücke}
Eine gewöhnliche kritische SU(4)-Kette hat links- und rechtslaufende Sektoren, mit $c_L=c_R=3$. Die übliche Kennzahl heißt dabei $c=3$, nicht sechs. Die chirale Differenz ist null. Eine rein chirale $E_8$-Naht hat dagegen Differenz acht. Ein zusätzliches ebenfalls nichtchirales $D_5$-Modell ändert daran nichts.

Eine mögliche physikalische Architektur benutzt einen zweidimensionalen Bulk, in dem gegenläufige innere Moden gegappt werden und chirale Randmoden bleiben. Aber ein solcher 2+1D-Bulk ist noch keine 3+1D-Raumzeit. Auch ein holomorpher Rand bedeutet nicht, dass der passende Bulk „leer“ ist: Eine gapped invertierbare Phase kann eine nichttriviale topologische Antwort tragen.

\section{Zwei Vierteldrehungen und der Index}
Die innere Glue-$\Z_4$ wirkt auf den 248 Strömen mit Phasen $1,i,-1,-i$ auf Sektoren der Größen $60,64,60,64$. Ihre Spur ist null. Die geometrische Vierteldrehung $e^{i\pi L_0/2}$ hat auf Stufe eins dagegen Spur $248i$. Die Operatoren sind somit nicht identisch. Eine Clock kann eine geometrische Drehung mit innerem Twist kombinieren; die beiden Faktoren bleiben getrennt. \src{S08, §8}

Die Quellen verbinden außerdem den Gitterindex $|\det A|$ mit dem Jones-Index eines passenden Gitterunternetzes. Ein logarithmischer Index kann eine Entropie- oder Inklusionsgröße liefern. Er legt für sich keinen physikalischen Hamiltonoperator fest. Die pauschale Gleichsetzung mit „der Entropie“ benötigt den konkreten Zustand und die verwendete Entropiedefinition. Ebenso ist „Normzeit ist RG-Zeit“ zunächst eine Zuordnung verschiedener mathematischer Rollen, kein allgemeiner Dynamiksatz. Die berichtete Untergitterzählfunktion $\prod_{k=0}^7\zeta(s-k)$ betrifft eine arithmetische Zählung und keine bereits hergeleitete physische Zeit.

\par\noindent\textbf{Arithmetischer Anschluss:} Kapitel~\ref{ch:shadow-arithmetic} prüft den Vergleich dieser Naht mit dem Primon-/Bost--Connes-Modell. Es unterscheidet die Zustandssummen und erklärt, weshalb endlich viele feste additive Takte nicht alle Primlogarithmen liefern.
